%==============================================================
%  lecons/analyse/bolzano.tex — Théorème de Bolzano-Weierstrass
%==============================================================

\lecontitle{Théorème de Bolzano--Weierstrass}{Analyse réelle — Suites et compacité}

%=============================================================
%  § 1 — ÉNONCÉ
%=============================================================
\secnum{navyblue}{1}{Énoncé et signification}

\begin{tcolorbox}[thmbox]
  \textbf{Théorème.}\enspace\index{Bolzano-Weierstrass (théorème de)}
  \textit{Toute suite bornée\index{suite bornée} de réels admet au moins une
  sous-suite\index{sous-suite} convergente.}

  \medskip
  Précisément : si $\exists\, M>0$ tel que $|u_n| \leq M$ pour tout $n$,
  alors il existe $\varphi : \mathbb{N} \hookrightarrow \mathbb{N}$
  strictement croissante et $\ell \in [-M,M]$ avec
  \[
    u_{\varphi(n)} \;\longrightarrow\; \ell \quad (n \to +\infty).
  \]
\end{tcolorbox}

\begin{significationintuitive}
Une suite bornée vit dans un domaine borné : elle revient infiniment souvent au
voisinage d'au moins un point — sa \emph{valeur d'adhérence}\index{valeur d'adhérence}.
On peut toujours en extraire un fil convergent. Ce résultat fonde la
\textbf{compacité séquentielle}\index{compacité séquentielle} des intervalles
fermés bornés de $\mathbb{R}$ et est indispensable pour les théorèmes
d'existence en analyse.
\end{significationintuitive}

\decsep{}

%=============================================================
%  § 2 — DÉMONSTRATION
%=============================================================
\secnum{navyblue}{2}{Démonstration}

\begin{ideepreuve}
On coupe $[-M,M]$ en deux ; la moitié contenant une infinité de termes est
gardée, et on y sélectionne un terme. En répétant, les intervalles emboîtés
se resserrent vers un point $\ell$, limite de la sous-suite construite.
\end{ideepreuve}

\rubriq{Preuve}

\begin{mdframed}[style=proofbar]
  \textit{Initialisation.} Posons $I_0 = [-M,M]$ : tous les termes de
  $(u_n)$ appartiennent à $I_0$, qui en contient donc une infinité. On pose
  $\varphi(0) = 0$, de sorte que $u_{\varphi(0)} \in I_0$ ; c'est le rang de
  base de l'extractrice.

  \medskip
  \textit{Induction.} Supposons construits, pour tout $j$ tel que
  $0 \leq j \leq k$ :
  \begin{itemize}
    \item un intervalle fermé $I_j$ de longueur $2M/2^j$ contenant une
          infinité de termes de $(u_n)$, avec $I_j \subset I_{j-1}$ dès que
          $j \geq 1$ ;
    \item un indice $\varphi(j)$ tel que $u_{\varphi(j)} \in I_j$, avec
          $\varphi(j) > \varphi(j-1)$ dès que $j \geq 1$.
  \end{itemize}
  On divise $I_k = [a_k, b_k]$ en ses deux moitiés et on choisit $I_{k+1}$
  parmi celles qui contiennent une infinité de termes. On extrait
  $\varphi(k+1) > \varphi(k)$ tel que $u_{\varphi(k+1)} \in I_{k+1}$.

  \medskip
  \textit{Convergence.} Les suites $(a_k)$ et $(b_k)$ sont adjacentes :
  $(a_k)$ croît et $(b_k)$ décroît (car $I_{k+1} \subset I_k$), tandis que
  $b_k - a_k = 2M/2^k \to 0$ ; elles convergent donc vers une même limite
  $\ell$.
  Par encadrement :
  \[
    a_k \;\leq\; u_{\varphi(k)} \;\leq\; b_k
    \quad\Longrightarrow\quad
    u_{\varphi(k)} \;\xrightarrow{k \to \infty}\; \ell.
  \]
  \hfill$\blacksquare$
\end{mdframed}

\rubriq{Reformulation : l'ensemble des valeurs d'adhérence}

\begin{mdframed}[style=proofbar]
  L'ensemble $E$ des valeurs d'adhérence d'une suite bornée est non vide
  (par Bolzano--Weierstrass), fermé (limite de valeurs d'adhérence) et borné.
  On définit :
  \[
    \limsup_{n\to\infty} u_n = \max E, \qquad
    \liminf_{n\to\infty} u_n = \min E.
  \]
  La suite converge si et seulement si $\limsup u_n = \liminf u_n$.
  \hfill$\blacksquare$
\end{mdframed}

\decsep{}

%=============================================================
%  § 3 — EXEMPLES, CONTRE-EXEMPLES, EXERCICES
%=============================================================
\secnum{exgreen}{3}{Exemples, contre-exemples et exercices}

\rubriq{Exemples}

\begin{mdframed}[style=exbar]
  \textbf{1.}\enspace%
  $u_n = \sin(n)$ est bornée par $1$ ; tout point de $[-1,1]$ en est une
  valeur d'adhérence.

  \smallskip\noindent
  \textbf{2.}\enspace%
  $u_n = {(-1)}^n\!\left(1 - \tfrac{1}{n+1}\right)$ admet $1$ et $-1$ comme
  valeurs d'adhérence (sous-suites paire et impaire).

  \smallskip\noindent
  \textbf{3.}\enspace%
  $v_n = \{n\sqrt{2}\}$ (partie fractionnaire de $n\sqrt{2}$) est bornée
  dans $[0,1)$ et dense dans $[0,1]$ :
  chaque point de $[0,1]$ est limite d'une sous-suite.
\end{mdframed}

\rubriq{Contre-exemples}

\begin{mdframed}[style=cexbar]
  \textbf{Borne nécessaire.}\enspace%
  $u_n = n$ est non bornée et n'admet aucune sous-suite convergente dans
  $\mathbb{R}$.

  \smallskip\noindent
  \textbf{Dimension infinie.}\enspace%
  Dans $\ell^2(\mathbb{N})$, les vecteurs de base $(e_n)$ vérifient
  $\|e_n\| = 1$ et $\|e_n - e_m\| = \sqrt{2}$ pour $n \neq m$ : aucune
  sous-suite n'est de Cauchy. La compacité séquentielle est spécifique à
  la dimension finie.
\end{mdframed}

\vspace{6pt}
\exolabel

\begin{mdframed}[style=exobar]
  \begin{enumerate}
    \item Montrer que $u_n = \cos(n\pi/6)$ est bornée et déterminer toutes
          ses valeurs d'adhérence.

    \item Soit $(u_n)$ bornée. Montrer que l'ensemble $E$ de ses valeurs
          d'adhérence est non vide, fermé, borné, et définir
          $\limsup u_n$ et $\liminf u_n$.

    \item Si $(u_n)$ est bornée et que toutes ses sous-suites convergentes
          tendent vers $\ell$, montrer que $(u_n) \to \ell$.

    \item Énoncer et démontrer Bolzano--Weierstrass dans $\mathbb{R}^d$
          par extraction successive sur chaque coordonnée.

    \item (Théorème de Heine) Montrer qu'une fonction continue sur un
          segment $[a,b]$ y est uniformément continue, en raisonnant
          par l'absurde à l'aide de Bolzano--Weierstrass.
  \end{enumerate}
\end{mdframed}

\leconfooter{B.\ Bolzano (1817) \& K.\ Weierstrass (1874)}
